\section{Background}\label{sec:background}

This section presents the basic Information Retrieval (IR) concepts used in the
project. Before applying the system to the Cranfield collection, we illustrate
the main steps using a toy corpus from the assignment. The toy corpus contains
three documents:

\begin{itemize}
    \item $d_1$: The star in our solar system provides heat and light.
    \item $d_2$: That Hollywood star walked the red carpet for the movie premiere.
    \item $d_3$: Astronomers observe distant stars and galaxies using telescopes.
\end{itemize}

\subsection{Preprocessing and Inverted Index}

The first stage in an IR system is preprocessing. In this project,
preprocessing includes tokenization, stopword removal, and inflection
reduction. Tokenization splits text into individual terms. After lowercasing
and tokenizing, the toy documents become:

\begin{itemize}
    \item $d_1$: ["the", "star", "in", "our", "solar", "system", "provides", "heat", "and", "light"]
    \item $d_2$: ["that", "hollywood", "star", "walked", "the", "red", "carpet", "for", "the", "movie", "premiere"]
    \item $d_3$: ["astronomers", "observe", "distant", "stars", "and", "galaxies", "using", "telescopes"]
\end{itemize}

Using the stopword set $\{$the, in, our, and, that, for$\}$, the processed
documents after stopword removal become:

\begin{itemize}
    \item $d_1$: ["star", "solar", "system", "provides", "heat", "light"]
    \item $d_2$: ["hollywood", "star", "walked", "red", "carpet", "movie", "premiere"]
    \item $d_3$: ["astronomers", "observe", "distant", "stars", "galaxies", "using", "telescopes"]
\end{itemize}

Stemming or lemmatization then reduces related word forms into a common
representation. For example, ``provides'' becomes ``provide'', ``walked''
becomes ``walk'', ``stars'' becomes ``star'', ``galaxies'' becomes
``galaxy'', and ``telescopes'' becomes ``telescope''. The final processed
documents are:

\begin{itemize}
    \item $d_1$: ["star", "solar", "system", "provide", "heat", "light"]
    \item $d_2$: ["hollywood", "star", "walk", "red", "carpet", "movie", "premiere"]
    \item $d_3$: ["astronomer", "observe", "distant", "star", "galaxy", "use", "telescope"]
\end{itemize}

An inverted index maps each term to the documents in which it occurs. For the
processed toy collection, the inverted index is shown in
Table~\ref{tab:toy_index}.

\begin{table}[h]
\centering
\caption{Inverted index for the processed toy corpus.}
\label{tab:toy_index}
\begin{tabular}{|l|l|}
\hline
\textbf{Term} & \textbf{Posting list} \\ 
\hline
star & $d_1, d_2, d_3$ \\ 
\hline
solar & $d_1$ \\ 
\hline
system & $d_1$ \\ 
\hline
provide & $d_1$ \\ 
\hline
heat & $d_1$ \\ 
\hline
light & $d_1$ \\ 
\hline
hollywood & $d_2$ \\ 
\hline
walk & $d_2$ \\ 
\hline
red & $d_2$ \\ 
\hline
carpet & $d_2$ \\ 
\hline
movie & $d_2$ \\ 
\hline
premiere & $d_2$ \\ 
\hline
astronomer & $d_3$ \\ 
\hline
observe & $d_3$ \\ 
\hline
distant & $d_3$ \\ 
\hline
galaxy & $d_3$ \\ 
\hline
use & $d_3$ \\ 
\hline
telescope & $d_3$ \\ 
\hline
\end{tabular}
\end{table}

\subsection{TF-IDF Representation}

The Vector Space Model (VSM) represents each document as a vector of term
weights. The most common weighting scheme is TF-IDF. Term Frequency (TF)
measures how often a term occurs in a document. Let $f_{t,d}$ denote the
frequency of term $t$ in document $d$. Inverse Document Frequency (IDF)
assigns higher weight to terms that occur in fewer documents:

\begin{equation}
IDF(t)=\log_{10}\left(\frac{N}{df_t}\right),
\end{equation}

where $N$ is the total number of documents and $df_t$ is the number of
documents containing term $t$. The final TF-IDF weight is computed as:

\begin{equation}
TF\text{-}IDF(t,d)=TF(t,d)\times IDF(t).
\end{equation}

For the toy collection, $N=3$. The term ``star'' appears in all three
documents, so its IDF value is:

\[
IDF(\text{star})=\log_{10}\left(\frac{3}{3}\right)=0.
\]

Therefore, although ``star'' occurs frequently, it receives zero TF-IDF
weight because it appears in every document and provides no discriminative
power for ranking. Terms that appear in only one document have:

\[
IDF=\log_{10}\left(\frac{3}{1}\right)=0.4771.
\]

\begin{table}[h]
\centering
\caption{TF-IDF weights in the toy corpus.}
\label{tab:toy_tfidf}
\begin{tabular}{|l|c|c|c|c|c|c|c|}
\hline
\textbf{Term} & \textbf{TF(d1)} & \textbf{TF(d2)} & \textbf{TF(d3)} & \textbf{IDF} & \textbf{TF-IDF(d1)} & \textbf{TF-IDF(d2)} & \textbf{TF-IDF(d3)} \\ 
\hline
star & 1 & 1 & 1 & 0.0000 & 0.0000 & 0.0000 & 0.0000 \\ 
\hline
solar & 1 & 0 & 0 & 0.4771 & 0.4771 & 0.0000 & 0.0000 \\ 
\hline
light & 1 & 0 & 0 & 0.4771 & 0.4771 & 0.0000 & 0.0000 \\ 
\hline
movie & 0 & 1 & 0 & 0.4771 & 0.0000 & 0.4771 & 0.0000 \\ 
\hline
premiere & 0 & 1 & 0 & 0.4771 & 0.0000 & 0.4771 & 0.0000 \\ 
\hline
galaxy & 0 & 0 & 1 & 0.4771 & 0.0000 & 0.0000 & 0.4771 \\ 
\hline
telescope & 0 & 0 & 1 & 0.4771 & 0.0000 & 0.0000 & 0.4771 \\ 
\hline
system & 1 & 0 & 0 & 0.4771 & 0.4771 & 0.0000 & 0.0000 \\ 
\hline
provide & 1 & 0 & 0 & 0.4771 & 0.4771 & 0.0000 & 0.0000 \\ 
\hline
heat & 1 & 0 & 0 & 0.4771 & 0.4771 & 0.0000 & 0.0000 \\ 
\hline
walk & 0 & 1 & 0 & 0.4771 & 0.0000 & 0.4771 & 0.0000 \\ 
\hline
red & 0 & 1 & 0 & 0.4771 & 0.0000 & 0.4771 & 0.0000 \\ 
\hline
carpet & 0 & 1 & 0 & 0.4771 & 0.0000 & 0.4771 & 0.0000 \\ 
\hline
astronomer & 0 & 0 & 1 & 0.4771 & 0.0000 & 0.0000 & 0.4771 \\ 
\hline
observe & 0 & 0 & 1 & 0.4771 & 0.0000 & 0.0000 & 0.4771 \\ 
\hline
distant & 0 & 0 & 1 & 0.4771 & 0.0000 & 0.0000 & 0.4771 \\ 
\hline
use & 0 & 0 & 1 & 0.4771 & 0.0000 & 0.0000 & 0.4771 \\ 
\hline
\end{tabular}
\end{table}

\subsection{Boolean Retrieval}

In Boolean retrieval, a query is matched using the inverted index. For the
query ``star light'', the term ``star'' occurs in $d_1$, $d_2$, and $d_3$,
while ``light'' occurs only in $d_1$.

If the system uses OR semantics, all three documents match at least one query
term:

\[
d_1, d_2, d_3.
\]

If the system uses AND semantics, only $d_1$ matches both query terms:

\[
d_1.
\]

\subsection{Cosine Similarity and Ranking}

In the Vector Space Model, both documents and queries are represented as
TF-IDF vectors. Consider the query:

\[
q = \text{``star light''}.
\]

The TF-IDF representation of the query is shown in
Table~\ref{tab:query_tfidf}.

\begin{table}[h]
\centering
\caption{TF-IDF representation of the query ''star light''.}
\label{tab:query_tfidf}
\begin{tabular}{|l|c|c|c|}
\hline
\textbf{Term} & \textbf{TF(Query)} & \textbf{IDF} & \textbf{TF-IDF(Query)} \\ 
\hline
star & 1 & 0.0000 & 0.0000 \\ 
\hline
light & 1 & 0.4771 & 0.4771 \\ 
\hline
\end{tabular}
\end{table}

The cosine similarity between a query vector $\vec{q}$ and a document vector
$\vec{d}$ is computed as:

\begin{equation}
\cos(q,d)=
\frac{\vec{q}\cdot\vec{d}}
{\|\vec{q}\|\|\vec{d}\|}.
\end{equation}

For document $d_1$:

\[
\vec{q}\cdot\vec{d_1}
=
(0.4771)(0.4771)
=
0.2276.
\]

The magnitude of the query vector is:

\[
||q||=\sqrt{0.4771^2}=0.4771.
\]

The magnitude of document $d_1$ is:

\[
||d_1||
=
\sqrt{
5(0.4771^2)
}
=
1.0668.
\]

Thus, the cosine similarity between the query and $d_1$ is:

\[
\cos(q,d_1)
=
\frac{0.2276}
{(0.4771)(1.0668)}
=
0.4472.
\]

Documents $d_2$ and $d_3$ do not contain the term ``light'', so their cosine
similarity with the query is zero.

Therefore, the ranking is:

\[
d_1 > d_2 = d_3.
\]

This ranking is desirable because the query ``star light'' refers to the
solar-system meaning of ``star'', and $d_1$ discusses heat and light from the
sun.

\subsection{Word Sense Ambiguity}

Consider the query ``movie star". The ideal document to retrieve is $d_2$
because it discusses a Hollywood actor and a movie premiere, which matches
the entertainment-related meaning of the query.

The TF-IDF retrieval system successfully retrieves $d_2$ because the term
``movie'' appears only in that document and therefore has a high IDF weight.
This makes $d_2$ more relevant than the other documents.

However, the word ``star'' introduces word sense ambiguity. In document
$d_1$, ``star'' refers to a celestial object in the solar system, while in
$d_2$ it refers to a famous actor. The Vector Space Model treats both meanings
as the same term because it relies only on lexical matching and does not
understand semantic context.

If the query were only ``star'', then all three documents would receive weak
or similar scores because the term appears in every document and therefore has
IDF zero in this toy collection. As a result, the system cannot effectively
distinguish between the different meanings of the word ``star''.
